\section{F. 이기적인 별자리}

\begin{frame} % No title at first slide
    \sectiontitle{F}{이기적인 별자리}
    \sectionmeta{
        출제자 -- 노태윤({\href{https://solved.ac/profile/areian13}{@areian13}})\\
        태그: \consolas {\#set, \#mst, \#greedy}\\
        예상 난이도 -- \textbf{\color{acG2}Gold2}
    }
    \begin{itemize}
        \item 제출 10번, 정답 0명 (정답률 0\%)
        \item 처음 푼 사람: -
    \end{itemize}
\end{frame}
%----------------------------------------------------------------------------
\begin{frame}{\textbf{F}. 이기적인 별자리}
    \begin{itemize}
        \item 일단 좌표가 중복되는 점에 대한 처리부터 해봅시다.
        \item 가장 간단한 방법은 모든 점을 표현할 수 있는 2차원 배열을 만드는 겁니다.\\
        이런! 좌표의 범위는 $-10\,000 \sim 10\,000$이니, $20\,001 \times 20\,001 \approx$ 4억짜리 배열을 만드는 건 무리일 듯 하네요..
        \item 대신 맵을 사용하면 됩니다.
        별의 개수는 2,000개이므로, 해시맵이든 트리맵이든 크게 문제 없습니다.
        \item 해시맵과 트리맵 모두 편의상 map이라 쓰겠습니다.\\
        $\texttt{map<int, map<int, int>{}> pos}$를 선언하여\\
        $\texttt{pos[x][y] = max(pos[x][y], b)}$로 갱신해주면 됩니다.
    \end{itemize}
\end{frame}

%----------------------------------------------------------------------------
\begin{frame}{\textbf{F}. 이기적인 별자리}
    \begin{itemize}
        \item $\texttt{map<pair<int, int>, int> pos}$를 쓰는 방법도 있습니다.\\
        $\texttt{pos[\{x, y\}] = max(pos[\{x, y\}], b)}$ 이런 식으로 쓰겠군요.

        \item 해시맵을 쓰기 위해선 키 타입에 대해 hash 함수와 equal 연산자(==)가 정의 돼있어야 합니다.
        하지만 C++에선 pair의 해시함수가 정의 돼있지 않습니다. 따라서 pair의 해시함수를 직접 정의해주셔야 합니다.\\
        언어별 정의 방법은 구글링 해보는 것을 추천합니다

        \item 트리맵은 키 타입에 대해 less 연산자(<)만 정의 돼있으면 됩니다.\\
        C++는 pair에 대해 less 연산자가 정의돼 있으므로 별도의 구현 없이 바로 사용할 수 있습니다.
    \end{itemize}
\end{frame}
%----------------------------------------------------------------------------
\begin{frame}{\textbf{F}. 이기적인 별자리}
    \begin{itemize}
        \item 이제 별을 이어봅시다.\\
        만약 중심별이란 개념이 없었다면, 이 문제는 매우 웰노운한 MST였습니다.\\
        $n$개의 점에 대해서 $O(n^2)$으로 간선들을 만드로 MST를 돌리면 되죠.
        \item 허나, 중심별에 의해 무작정 별을 이을 수 없습니다.\\
        MST를 진행하면서 어떤 별 $u, v$를 연결하고자 할 때,\\
        $u, v$가 각자 속한 별자리(별 집합)에 이미 중심별이 있는가를 판단하고,\\
        양 쪽 모두 중심별이 있다면 연결해선 안됩니다.
        \item 이를 구현하는 방법은 크게 두가지가 있습니다.
    \end{itemize}
\end{frame}

%----------------------------------------------------------------------------
\begin{frame}{\textbf{F}. 이기적인 별자리}
    \begin{itemize}
        \item 첫번째 방법입니다.\\
        일단 MST 과정에서 $u, v$ 모두 비중심별이라면, 고민할 것 없이 $\texttt{union}$해줍니다.
        \item 만약 $u$는 중심별, $v$는 비중심별이라면, $u$가 $v$의 $\texttt{parent}$가 되도록 분리 집합을 형성합니다.\\
        즉, $\texttt{find}(v)==u$가 되도록 구현하면 됩니다.
        \item 이후 다른 간선의 별 $u, v$를 $\texttt{union}$하기 전,\\
        $\texttt{find}(u)$와 $\texttt{find}(v)$가 모두 중심별이라면 $\texttt{union}$하지 않고 넘어가면 됩니다.
        \item 이렇게 MST를 진행하면\\
        각각의 분리 집합이 하나의 중심별을 $\texttt{root}$ 노드로 갖는 최소 비용 집합이 됩니다.
    \end{itemize}
\end{frame}
%----------------------------------------------------------------------------
\begin{frame}{\textbf{F}. 이기적인 별자리}
    \begin{itemize}
        \item 첫번째 방법은 직관적이나, 루트 노드가 무엇인지, 비/중심별을 판별하는 등의 구현상의 귀찮음이 많습니다.
        두번째 방식은 매우 그리디합니다.
        \item 결론적으로, 간선을 생성하는 단계에서 만약 두 별 $u, v$가 모두 중심별이라면, 간선으 비용을 0으로 설정하면 됩니다.\\
        이렇게 하고 간선들을 정렬하고 MST를 돌리면,
        비용이 0인 중심별들 사이의 간선이 최우선으로 연결됩니다.
        \item 이후 이미 어떤 중심별과 연결된 비중심별 $u,v$가 $\texttt{union}$을 시도할 때,
        개념적으로는 본인들이 속한 별자리가 별개의 별자리일지라도, 이미 중심별들끼리 하나의 집합을 이루고 있으므로 $\texttt{union}$되지 않습니다.
    \end{itemize}
\end{frame}
%----------------------------------------------------------------------------
\begin{frame}{\textbf{F}. 이기적인 별자리}
    \begin{itemize}
        \item 아이디어가 중요한 문제였고, 시간 복잡도 자체는 일반적인 MST와 다를 바 없기 때문에\\

        \bigskip

        간선을 만드는 비용 $O(n^2)=O(e)$\\
        정렬하는 비용 $O(e \log e)$\\
        MST를 돌리는 비용 $O(e \log n)$

        \bigskip

        입니다.
    \end{itemize}
\end{frame}